\documentclass[11pt,a4paper]{article}
\usepackage{nesarticle}
\begin{document}
\heading{Homework 1 · Problem solutions}{Instructor guide · Lectures 1 and 2}
\section*{Problem 1. Two tickets, one weekend}
\paragraph{(a) Rankings and rationality.} All three programmes have utility 3: the score pairs are $(6,3)$, $(3,9)$ and $(3,4)$. Thus they are mutually indifferent. A real-valued representation gives completeness and transitivity. The minimum of two continuous linear functions is continuous, so the represented weak-preference graph is closed.
\paragraph{(b) Geometry.} The level-6 indifference set is the polyline connecting
\[
 (0,6),\qquad(12/5,6/5),\qquad(6,0).
\]
The weak upper contour set is
\[
 \{x\ge0:2x_1+x_2\ge6,\quad x_1+3x_2\ge6\}.
\]
For $0\le x_1\le12/5$, the indifference set follows $x_2=6-2x_1$. Mira's score is the smaller one there (strictly smaller before the kink). For $12/5\le x_1\le6$, it follows $x_2=(6-x_1)/3$, and Anton's score is the smaller one. At the kink both scores are 6. The upper contour is above both lines, not a union of the two half-spaces.
\paragraph{(c) Properties.} At any reference utility level $k\ge0$, the upper contour is the intersection of $X$ with two half-spaces $2x_1+x_2\ge k$ and $x_1+3x_2\ge k$, so preferences are convex. They are not strictly convex: $x=(0,4)$ and $y=(1,2)$ both have utility 4, and their midpoint $(1/2,3)$ also has utility 4.

Both linear scores have strictly positive coefficients. If a bundle increases coordinatewise with at least one strict improvement, \emph{both} scores strictly rise. More precisely, both score increases have a strictly positive minimum, so their minimum also rises strictly. This proves strong monotonicity. Increasing either coordinate by an arbitrarily small positive amount then proves LNS throughout the orthant.
\paragraph{(d) Displayed scores.} Since $u(X)=[0,\infty)$, squaring is strictly increasing on the attained range and preserves every comparison. But $v$ is not concave. On $x_2=0$, $u(x)=x_1$ and $v(x)=x_1^2$. For example,
\[
 v(1,0)=1<\tfrac12v(0,0)+\tfrac12v(2,0)=2.
\]
Changing the numerical scale need not preserve concavity.
\paragraph{(e) Convex representation.} Set $x=(0,6)$ and $y=(6,0)$. Both have utility 6, while their midpoint $(3,3)$ has utility 9. Any representation $g$ must satisfy $g(x)=g(y)<g((x+y)/2)$, whereas convexity would require $g((x+y)/2)\le[g(x)+g(y)]/2=g(x)$. Thus no convex representing function exists. Convexity of preferences does not imply convexity of a utility index.
\clearpage
\section*{Problem 2. One more lesson, one less coffee}
\paragraph{(a) Existence, exhaustion and uniqueness.} The budget is nonempty and compact for $p>0,m>0$, and utility is continuous, so an optimum exists. Utility is strictly increasing in $a$ and $q$, hence an optimum exhausts the budget. Substitute $a=m-pq$ to obtain
\[
 \max_{0\le q\le m/p}\{m-pq+6\log(1+q)\}.
\]
The second derivative with respect to $q$ is $-6/(1+q)^2<0$. The reduced problem has a unique maximiser; budget exhaustion then uniquely determines $a$. The original utility is not strictly concave in the pair $(a,q)$ because it is affine in $a$; do not use that false premise to justify uniqueness.
\paragraph{(b) Kuhn--Tucker system.} With constraints written as nonnegative,
\[
 \mathcal L=a+6\log(1+q)+\lambda(m-a-pq)+\mu_a a+\mu_q q.
\]
Stationarity:
\[
 1-\lambda+\mu_a=0,\qquad
 \frac6{1+q}-\lambda p+\mu_q=0.
\]
Primal feasibility: $a,q\ge0$, $m-a-pq\ge0$. Dual feasibility: $\lambda,\mu_a,\mu_q\ge0$. Complementary slackness:
\[
 \lambda(m-a-pq)=0,\qquad \mu_a a=0,\qquad\mu_q q=0.
\]
The objective is concave and differentiable on a neighbourhood of the consumption domain; the constraints are affine. KKT candidates satisfying these conditions are globally optimal. The reduced problem in (a) already establishes uniqueness.
\paragraph{(c) Demand.} The reduced derivative is $-p+6/(1+q)$, strictly decreasing in $q$. Thus
\[
 q^*(p,m)=\min\left\{\frac m p,\ \max\left\{\frac6p-1,0\right\}\right\},\qquad
 a^*(p,m)=m-pq^*(p,m).
\]
The full classification is
\[
 (a^*,q^*)=
 \begin{cases}
 (m,0),&p\ge6,\\[1mm]
 (0,m/p),&0<p<6,\quad 0<m\le6-p,\\[1mm]
 (m-6+p,6/p-1),&0<p<6,\quad m>6-p.
 \end{cases}
\]
At $p=6$, demand for tutoring is still uniquely zero: the reduced derivative is zero only at zero and negative thereafter. At $m=6-p$ with $p<6$, $a=0$ and $q=6/p-1=m/p$; this is the boundary between the two positive-tutoring regimes, not a multiplicity of optima.

The corresponding multipliers are
\[
 \lambda=
 \begin{cases}
 1,&p\ge6,\\
 6/(p+m),&0<p<6,\quad m\le6-p,\\
 1,&0<p<6,\quad m>6-p,
 \end{cases}\qquad
 \mu_a=\lambda-1.
\]
When $q>0$, $\mu_q=0$; when $q=0$, $\mu_q=p-6\ge0$. These formulas agree on the boundaries. In particular, a binding zero lower bound on a good need not have a strictly positive multiplier.
\paragraph{(d) Two incomes.} At $p=2$,
\[
 \begin{array}{c|ccc}
 m&a^*&q^*&\lambda\\\hline
 2&0&1&3/2\\
 8&4&2&1
 \end{array}
\]
With $m=2$, the no-borrowing/nonnegativity constraint on $a$ binds; the student would prefer to buy more tutoring at these prices but cannot finance it. For $m<4$, each additional unit of income buys $1/2$ more tutoring hour. With $m=8$, tutoring is at its unconstrained preferred level of 2 hours and additional income goes to $a$. Quasilinearity does not eliminate low-income corner effects.
\paragraph{(e) Indirect utility and Roy.} In the open region $0<p<6$, $m>6-p$, both goods are consumed and
\[
 v(p,m)=m-6+p+6\log(6/p),\qquad
 v_m=1,\qquad v_p=1-\frac6p.
\]
Hence $-v_p/v_m=6/p-1=q^*$. The derivative $v_m$ is the marginal value of income in the chosen utility units, equal to the budget multiplier. In this region one extra unit of income buys one more unit of $a$ and raises utility by 1. The multiplier is not invariant to rescaling the utility representation.

\end{document}
